\documentclass[11pt]{article}
\usepackage{coursenotes}
\begin{document}
\notesheader{9}{Generative AI in Causal Analysis}{Learned simulators, and variables generated by language models}

\begin{decision}
Two new uses of generative AI arrive in every analytics team. One builds a simulator of the business to test pricing or
targeting pipelines before deploying them. The other uses a language model to turn reviews, chats, or listings into variables
for an analysis. Both can improve a causal analysis, or bias it without anyone noticing. In both cases the question is the one this course
always asks: where does the truth come from, and how is it checked?
\end{decision}

\begin{goals}
  \item show that a learned simulator is a causal model resting on the same assumptions as the g-formula;
  \item plant effects and confounding, and use a simulator as an estimator testbed and as a pipeline testbed;
  \item decompose a policy's regret and explain home advantage, backtests, and support;
  \item distinguish planted truth from asserted truth (silicon samples, CEVAE);
  \item derive the bias from a model-coded outcome, treatment, or modifier, and correct it with a gold subsample (PPI);
  \item use text embeddings as controls without leakage or overlap collapse.
\end{goals}

%=====================================================================================
\sectionone

\subsection*{Part A: Learned simulators}

\subsection{From standardization to a simulator}

The g-formula (Meeting~7) computes $\E[Y(d)]$ by averaging $\mu_d(x)$ over $X$. A \emph{learned simulator} extends this from averages
to distributions. It fits a generative model in pieces that follow the causal order, $\hat p(x)$, $\hat p(d \mid x)$, and
$\hat p(y \mid d, x)$, each of which can be any generative model (a parametric fit, a mixture, a normalizing flow, a variational
autoencoder, a boosted tree with a fitted noise distribution). Sampling $x$, then $d$, then $y$ yields synthetic logs. Sampling $y$ for a
\emph{chosen} $d$ yields potential outcomes, so the simulator can evaluate any policy, including budgeted, sequential, or dynamic
rules that have no simple estimator.

\begin{prop}{A simulator's effects are g-formula effects}{simg}
If $\hat p(x) = p(x)$ and $\hat p(y \mid d, x) = p(y \mid d, x)$, the simulator's average outcome under action $d$ is
$\sum_x \E[Y \mid D=d, X=x]\,p(x)$. This equals $\E[Y(d)]$ if and only if the g-formula does, that is, under consistency, conditional
exchangeability given $X$, and overlap.
\end{prop}
\begin{proof}
Setting $D = d$ and sampling $Y$ from $p(y \mid d, x)$ gives mean $\E[Y \mid D=d, X=x]$ at each $x$; averaging over $p(x)$ is the g-formula,
which equals $\E[Y(d)]$ under the three assumptions (Meeting~7).
\end{proof}

So a simulator fitted to logs reproduces the observed joint distribution faithfully and still gets effects wrong if an unmeasured
confounder drove treatment: the selection bias of Meeting~1 is built into $p(y \mid d, x)$. Its realism is evidence about the data,
not about causation.

\subsection{Planting the truth}

The analyst can \emph{decide} the truth. \emph{Planting an effect}: keep $\hat p(x)$ and a control-outcome model, and define treated
outcomes as control outcomes plus a chosen $\tau^{\text{plant}}(x)$, often calibrated so that the simulator reproduces a past
randomized test. \emph{Planting confounding}: add a latent $U$ that raises both the chance of treatment and the outcome, with chosen
strengths, and hide it from the analyst. Because covariates, assignment patterns, and noise come from real data, the synthetic data
have the business's awkward features (skew, rare segments, poor overlap) rather than a textbook $y = a + bx + \varepsilon$. That realism is
what persuades stakeholders.

\subsection{Two uses: estimator testbed and pipeline testbed}

\paragraph{Estimator testbed.} Generate many datasets, run an estimator on each, compare with the planted truth: bias, coverage,
behavior under poor overlap. This checks code and statistics on this business's data shape.

\paragraph{Pipeline testbed.} Businesses deploy pipelines (a model feeding an optimizer feeding a decision), not estimators.
\begin{enumerate}
  \item Generate synthetic logs, with planted effects and, if wanted, planted confounding.
  \item Run each candidate pipeline exactly as in production: fit, optimize, output a policy.
  \item Evaluate each policy on the simulator's potential outcomes, where its value is known.
  \item Compute the \emph{oracle} policy from the planted truth, and each pipeline's \emph{regret}: oracle value minus pipeline value.
\end{enumerate}

\begin{center}
\begin{tikzpicture}[>={Stealth[length=2mm]}, every node/.style={font=\small}, box/.style={draw=noteblue, fill=bluetint, rounded corners=1pt, minimum height=0.9cm, align=center, text width=2.1cm}]
  \node[box] (real) at (0,0) {Real logs and past tests};
  \node[box] (sim) at (3.0,0) {Simulator\\(planted truth)};
  \node[box] (syn) at (6.0,0) {Synthetic logs};
  \node[box] (pipe) at (9.0,0) {Pipeline:\\model $\to$ optimizer};
  \node[box] (pol) at (9.0,-1.8) {Policy};
  \node[box, draw=sbgreen, fill=greentint] (reg) at (3.0,-1.8) {Value on true potential outcomes; regret vs oracle};
  \draw[->] (real) -- node[above, font=\scriptsize]{fit} (sim);
  \draw[->] (sim) -- node[above, font=\scriptsize]{sample} (syn);
  \draw[->] (syn) -- (pipe); \draw[->] (pipe) -- (pol);
  \draw[->] (pol) -- (reg); \draw[->, dashed] (sim) -- node[left, font=\scriptsize]{truth} (reg);
\end{tikzpicture}

\small Figure 1. The pipeline testbed. The pipeline never sees the truth; only the evaluation does.
\end{center}

\begin{prop}{Regret counts only misclassified units}{regret}
Without a budget, the oracle treats $x$ if and only if its net value $v(x) > 0$. For any policy $\pi$,
\[
  V(\pi^\ast) - V(\pi) = \E\big[\,|v(X)|\;\ind{\pi(X) \ne \pi^\ast(X)}\big].
\]
\end{prop}
\begin{proof}
By the policy-value identity (Meeting~5), $V(\pi^\ast) - V(\pi) = \E[(\pi^\ast(X) - \pi(X))v(X)]$. Where the two agree the term is zero;
where $\pi^\ast = 1, \pi = 0$ it is $v > 0$; where $\pi^\ast = 0, \pi = 1$ it is $-v = |v|$.
\end{proof}

Mistakes near the break-even cost little; mistakes on units with large positive or negative net value cost a lot. A pipeline can
estimate effects poorly and still have low regret if its errors fall where $|v|$ is small. This is the reason decision-focused learning and policy learning (Meeting~6) train on the decision's value
rather than on estimation error: the weighted-classification result there is the sample version of this one.

\subsection{What makes a simulator convincing}

\begin{itemize}
  \item \textbf{Backtests against past interventions.} The simulator should reproduce, out of sample, the measured effects of past
  experiments or policy changes. This is the only check on its causal part.
  \item \textbf{Staying within support.} A simulator will produce outcomes for prices or rules never seen in the logs without warning.
  Keep evaluated policies within the training range; treat answers outside it as extrapolation.
  \item \textbf{Compounding error.} In sequential simulators small one-step errors accumulate; validate multi-step trajectories.
  \item \textbf{Home advantage.} A pipeline whose models share the simulator's model class tends to look good because it matches the
  data-generating process. Build simulators of more than one class and trust only rankings that hold across them.
\end{itemize}

\subsection{Planted truth and asserted truth}

In a simulator the analyst \emph{plants} the effect and the model supplies realism. In a \emph{silicon sample}, a language model is asked
how people would respond (``You are a 35-year-old parent in Shanghai. Would you order this at ¥39?'') and its answers are treated as
potential outcomes: the truth is \emph{asserted} by the model's priors, learned from text that may include the very studies being
replicated. Documented problems include too little variance, stereotyped subgroups, sensitivity to wording, and agreement with
published results that reflects memorization. Silicon samples can draft survey questions or set priors for power calculations; they
settle nothing causal unless validated out of sample against real experiments in a similar setting. The causal effect VAE (CEVAE),
which infers a hidden confounder from proxies, is on the asserted side too: its identification rests on structure it does not state
(Handout H2).

\subsection*{Part B: Variables generated by language models}

\subsection{Language models as measurement instruments}

A language model turns text into variables: a review's sentiment, whether a chat mentions a late delivery, an embedding of a menu.
Each is a measurement with error, and the error's consequence depends on the variable's role. For a binary true variable $Y$ and model
label $f$, write the false-positive rate $\alpha = \Pr(f=1 \mid Y=0)$ and false-negative rate $\beta = \Pr(f=0 \mid Y=1)$.

\begin{prop}{A model-coded outcome}{outcome}
If $(\alpha, \beta)$ are the same in both arms (non-differential error) and $\alpha + \beta < 1$, then
$\E[f \mid D=1] - \E[f \mid D=0] = (1 - \alpha - \beta)\{\E[Y \mid D=1] - \E[Y \mid D=0]\}$. If they differ by arm, $(\alpha_d, \beta_d)$,
\[
  \E[f \mid D=1] - \E[f \mid D=0] = (1-\alpha_1-\beta_1)p_1 - (1-\alpha_0-\beta_0)p_0 + (\alpha_1 - \alpha_0),
\]
with $p_d = \E[Y \mid D=d]$.
\end{prop}
\begin{proof}
$\E[f \mid D=d] = \alpha_d(1 - p_d) + (1 - \beta_d)p_d = \alpha_d + (1 - \alpha_d - \beta_d)p_d$; subtract across arms.
\end{proof}

Non-differential error shrinks the effect toward zero by the factor $1 - \alpha - \beta$. \emph{Differential} error, which arises whenever
the treatment changes the text in ways that affect the model's accuracy, biases the effect in either direction and can create an effect
from nothing.

\begin{prop}{A model-coded effect modifier}{modifier}
If a binary modifier $V$ splits 50/50 and the model's coded version $\tilde V$ has symmetric misclassification rate $q < \tfrac12$, then
$\tau(\tilde V = 1) - \tau(\tilde V = 0) = (1 - 2q)\{\tau(V = 1) - \tau(V = 0)\}$.
\end{prop}
\begin{proof}
With a 50/50 split, $\tau(\tilde V = 1) = (1-q)\tau_1 + q\tau_0$ and $\tau(\tilde V = 0) = q\tau_1 + (1-q)\tau_0$; subtract.
\end{proof}

A coded \emph{control} leaves residual confounding in proportion to its error, and a coded \emph{treatment} misclassifies exposure, usually
biasing toward zero. In every role, accuracy must be \emph{measured} against gold labels, separately by arm or group, and the coded field
must be built from text fixed before treatment.

\subsection{Correcting with a gold subsample}

Let a model label all $N$ units ($f_i$) and humans label a random subsample of $n$ ($Y_i$).

\begin{prop}{Prediction-powered inference (PPI) for a mean}{ppi}
\[
  \hat\theta_{\text{PP}} = \frac1N\sum_{i=1}^N f_i + \frac1n\sum_{i \in \text{gold}}(Y_i - f_i)
\]
is unbiased for $\E[Y]$ whatever the model's quality. If the gold sample is drawn independently of the $N$ units used for the first
term, $\Var(\hat\theta_{\text{PP}}) = \Var(f)/N + \Var(Y - f)/n$.
\end{prop}
\begin{proof}
The first term has mean $\E[f]$; the second, a random sample of $Y - f$, has mean $\E[Y] - \E[f]$; add. Independent terms' variances add.
\end{proof}

The second term, the \emph{rectifier}, removes the model's average error. A better model makes $Y - f$ less variable and the interval
narrower; a worse one widens it but never biases it. For an effect on a coded outcome, apply PPI in each arm, with gold labels drawn in
each arm: that is what catches differential error. Design-based supervised learning (DSL) extends the idea to regression coefficients and
to gold samples drawn with known unequal probabilities.

\subsection{Embeddings as controls}

Text often carries confounders that structured fields miss (a menu reveals a restaurant's positioning, which drives both its prices and its
orders). An \emph{embedding} can enter double machine learning as part of $X$. Two failure modes:
\begin{itemize}
  \item \textbf{Leakage.} Text written or edited after treatment (reviews posted after a price change) contains the treatment or its
  consequences: a bad control.
  \item \textbf{Overlap collapse.} A rich embedding can predict treatment almost perfectly, because text describes what drove treatment. The
  treatment residual vanishes and the estimate rests on a few units.
\end{itemize}
Use only text fixed before the decision, check overlap after adding the embedding, and compare estimates with and without it.

\begin{pitfall}[generative-AI errors]
Treating a simulator's realism as causal validity; ranking pipelines on one simulator class; reading a silicon sample as data; validating an
LLM label on pre-treatment text only; gold labels from one arm; embeddings of post-treatment text.
\end{pitfall}

%=====================================================================================
\sectiontwo

\subsection*{Part A: choosing a courier-bonus pipeline on two simulators}

\paragraph{Setting.} QuickBite pays bonuses to couriers for working evening peaks, to the top 30\% under a budget. Each extra peak hour is worth
¥30 and a bonus costs ¥60, so the net value of paying courier $x$ is $v(x) = 30\,\tau(x) - 60$. Two simulators were fitted to last year's
randomized bonus test (6{,}000 couriers): one from neural networks, one from boosted trees, each giving its own planted $\mu_0(x)$ and $\tau(x)$.
Each generates synthetic logs of 4{,}000 couriers in which past bonuses were targeted on courier features, with and without a planted hidden
confounder (motivation) that raises both bonus receipt and hours. Three pipelines: (1) \emph{outcome prediction}: pay those predicted to work
the most peak hours if paid; (2) a \emph{DR-learner} with cross-fitted boosting, paying the highest estimated effects; (3) a \emph{neural
T-learner}. All three are plug-in rules in the sense of Meeting~6: estimate, then rank and pay within the budget. Values are net yuan per 1{,}000 couriers, averaged over 20 synthetic datasets (\texttt{checks/m09\_sim.py}).
\begin{center}
\begin{tabular}{llcccc}
\toprule
Simulator & Hidden & Oracle & \multicolumn{3}{c}{Regret (¥ per 1{,}000 couriers)} \\
\cmidrule(lr){4-6}
 & confounder & value & Outcome & DR-learner & Neural T \\
\midrule
Neural  & none     & 22{,}761 & 18{,}838 & 1{,}891 & 858 \\
Neural  & moderate & 22{,}761 & 17{,}362 & 2{,}038 & 968 \\
Boosted & none     & 23{,}464 & 18{,}680 & 2{,}200 & 1{,}621 \\
Boosted & moderate & 23{,}464 & 17{,}083 & 2{,}383 & 1{,}668 \\
\bottomrule
\end{tabular}
\end{center}

\paragraph{Reading the table.}
\begin{itemize}
  \item Paying on predicted outcomes gives up about 80\% of the attainable value everywhere: it pays couriers who would work peaks anyway.
  \item The neural T-learner has the lowest regret on both simulators, but its lead over the DR-learner halves on the boosted simulator (from
  ¥1{,}033 to ¥579 per 1{,}000, SE about ¥100): part of its advantage on the neural simulator is home advantage. Because the ranking holds on
  both simulators, it deserves some trust.
  \item The planted confounder barely changes the causal pipelines' regret. It shifts everyone's estimated effect by a similar amount, which
  distorts levels but not the ranking, and a fixed-budget policy uses only the ranking (Meeting~5). A break-even threshold rule would suffer.
\end{itemize}

\subsection*{Part B: a model-coded outcome with a gold subsample}

\paragraph{Setting.} QuickBite randomized a ``Verified Kitchen'' badge across 20{,}000 restaurants (10{,}000 per arm). The outcome is whether a
restaurant's next review mentions a food-quality problem, coded by a language model. Humans labeled 1{,}000 random reviews in each arm.

\paragraph{Step 1: model labels only.} The model flags 8.0\% of treated and 10.0\% of control reviews: an effect of $-0.020$ (SE $0.004$,
interval $[-0.028, -0.012]$), apparently clear.

\paragraph{Step 2: the rectifiers.} On gold reviews, the mean of (human label minus model label) is $+0.012$ in the treated arm and $+0.002$ in the
control arm. The model misses more complaints about badged kitchens, whose reviewers write more politely: differential error.

\paragraph{Step 3: PPI.} Treated $0.080 + 0.012 = 0.092$; control $0.100 + 0.002 = 0.102$; effect $-0.010$. With $\widehat\Var(Y - f) \approx 0.05$ in each arm,
$\widehat\SE_{\text{treated}} = \sqrt{0.080 \times 0.920/10{,}000 + 0.05/1{,}000} = 0.0076$, $0.0077$ for control, so the effect's SE is $0.0108$ and its
interval $[-0.031, 0.011]$. Half the apparent effect was a measurement artifact, and the rest is not distinguishable from zero. Gold labels
alone would give SE $0.0132$; PPI's is about 18\% smaller.

\begin{exercise}[computation]
Suppose the model's errors were non-differential, with $\alpha = 0.02$ and $\beta = 0.20$ in both arms, and the model-coded effect were $-0.020$.
What is the effect on true complaints? What if instead $\alpha$ were $0.02$ in control and $0.01$ for badged kitchens, with $\beta = 0.20$ in both and
true complaint rates $0.10$ in both arms?
\end{exercise}
\answer{Non-differential: $-0.020/(1 - 0.02 - 0.20) = -0.026$; coded labels understate the effect. Differential with no true effect: the coded
difference is $(1 - 0.01 - 0.20)(0.10) - (1 - 0.02 - 0.20)(0.10) + (0.01 - 0.02) = 0.079 - 0.078 - 0.010 = -0.009$: a spurious reduction of
$0.009$ created entirely by the model's different false-positive rates.}

\begin{exercise}[judgment]
A consultant estimates QuickBite's price elasticity for a new dessert line by asking a language model to role-play 2{,}000 customer personas
at three prices, and reports $-1.8$ with a tight interval. What would you ask before using the number, and what could the personas legitimately
be used for?
\end{exercise}
\answer{Ask: has the persona method reproduced, out of sample, elasticities from QuickBite's own price tests in similar categories? Is the
interval's tightness just the model's low variance? Does the estimate move with prompt wording? Is $-1.8$ a textbook value the model recalls?
Without validation it is an asserted, not a measured, truth. Legitimate uses: drafting the menu test's design, choosing price points, setting a
prior for the power calculation of a real price test.}

%=====================================================================================
\sectionthree

\begin{extension}[Calibrated semi-synthetic benchmarks]
RealCause (Neal, Huang, and Raghupathi, 2020) and Credence (Parikh et al., 2022) fit deep generative models to real covariates and treatments and
let the user plant effects and confounding, with diagnostics for realism and their limits.
\end{extension}

\begin{extension}[Three routes on one testbed (optional lab or assignment)]
On the boosted simulator of Section~2, compare three pipelines over 20 replications by regret against the oracle: (1)
plug-in predict-then-optimize, a DR-learner thresholded at break-even; (2) a decision-focused learner, a gradient-boosted classifier
with label $\ind{\Gamma_i > 0}$ and weight $|\Gamma_i|$ from cross-fitted doubly robust scores (Meeting~6); (3) a depth-2 policy tree.
Repeat under a budget, replacing $\Gamma_i$ with $\Gamma_i - \lambda c_i$. Report where each pipeline's misclassified units sit relative
to break-even. Which route wins depends on where the effect model's errors fall, which is what the testbed can show.
\end{extension}

\begin{extension}[Foundation-model estimators]
In-context estimators trained on simulated causal worlds (Do-PFN) push the simulator idea to its limit: the estimator's assumptions live in the
prior over worlds. Handout H2 discusses them with neural estimators generally.
\end{extension}

\subsection*{Annotated reading}
\begin{readinglist}
\reading{Angelopoulos, Bates, Fannjiang, Jordan, and Zrnic (2023), ``Prediction-powered inference'', \emph{Science}}{PPI for means, regressions, and
beyond}{Main text}{2}
\reading{Egami, Hinck, Stewart, and Wei (2023), ``Using imperfect surrogates for downstream inference'', \emph{NeurIPS}}{Design-based supervised
learning for LLM-coded variables}{Sections 1--3}{3}
\reading{Veitch, Sridhar, and Blei (2020), ``Adapting text embeddings for causal inference'', \emph{UAI}}{Text as a confounder, with
embeddings}{Sections 1--3}{3}
\reading{Argyle et al.\ (2023), ``Out of one, many'', \emph{Political Analysis}; Horton (2023), ``Large language models as simulated economic
agents'' (NBER)}{The case for silicon samples, to read critically}{Main text}{1}
\reading{Parikh, Varjao, Xu, and Tchetgen Tchetgen (2022), ``Validating causal inference methods'', \emph{ICML}}{Credence: calibrated simulators
for method choice}{Sections 1--4}{3}
\reading{Melnychuk, Frauen, and Feuerriegel (2022), ``Causal transformer for estimating counterfactual outcomes'', \emph{ICML}}{Sequence models for
time-varying treatments}{Sections 1--4}{3}
\reading{Handout H2, \emph{Neural and foundation-model estimators}}{TARNet to Do-PFN, with what each assumes}{All}{2}
\end{readinglist}

\end{document}
